\section{Discussion}
\label{sec:discussion}

\subsection{Comparison with GPU Spectroscopy Codes}

The throughput of our batch forward model---$10{,}708$ spectra per second on a V100S---is
% Source: benchmark_summary.json/headline_numbers/batch_forward_max_throughput_spectra_per_sec = 10708.1
comparable to the $\sim 10{,}000$ spectra per second reported by \citet{Kawahara2022} for ExoJAX on the same GPU architecture.
% Source: e2e_pipeline_results.json/references/note
This is noteworthy because the two codes solve fundamentally different physics: ExoJAX computes molecular and atomic cross-sections for exoplanet atmospheric retrieval (a single-species, many-line opacity sum), while our code evaluates a multi-element Saha--Boltzmann chain with iterative charge balance, Boltzmann fitting, compositional closure, and Voigt broadening.
The comparable throughput suggests that the JAX/XLA compilation stack achieves similar GPU utilization efficiency across these different spectral workloads, and that the per-spectrum computational cost is dominated in both cases by the Voigt profile matrix---an $(N_\mathrm{wl}, N_\mathrm{lines})$ outer product that maps naturally to GPU parallelism.

The comparison with HELIOS-K \citep{Grimm2021} highlights a different tradeoff.
HELIOS-K is implemented in CUDA C++ with hand-tuned kernels optimized for specific GPU architectures, achieving higher raw throughput for opacity calculations than a JAX-based approach can match.
However, the JAX framework provides automatic differentiation, which enables gradient-based optimization (e.g., joint fitting of $T$, $n_e$, and $\bm{C}$ via L-BFGS-B) and Bayesian inference (via Hamiltonian Monte Carlo or variational methods) without manually deriving and implementing Jacobians.
For CF-LIBS applications, where the inversion problem is as important as the forward model, this capability is a significant practical advantage that justifies the modest performance overhead relative to hand-tuned CUDA.

\subsection{Anderson Acceleration: An Honest Assessment}

The Anderson acceleration results merit candid discussion.
The average iteration reduction of $1.6\times$ at tolerance $10^{-6}$ is modest,
% Source: benchmark_summary.json/headline_numbers/anderson_avg_iteration_reduction = 1.62
and the optimal depth of $m = 1$ at this tolerance means the algorithm reduces to a simple secant-like update rather than exploiting the full multi-step history.
% Source: benchmark_summary.json/headline_numbers/anderson_optimal_M = 1

The underlying reason is that the Saha--Boltzmann charge-balance iteration for typical LIBS plasmas ($T \sim 0.6$--\SI{1.5}{\electronvolt}, $n_e \sim 10^{15}$--$10^{17}$~\si{\per\cubic\centi\metre}) is already well-conditioned: the Picard contraction rate is $|g'(n_e^*)| \sim 0.3$--$0.7$, meaning simple iteration converges in 3--8 steps.
Anderson acceleration delivers its greatest benefit for problems with contraction rates near unity, where Picard iteration is slow or fails to converge.
For the charge-balance problem, such conditions arise at higher temperatures ($T > \SI{2}{\electronvolt}$) with significant double or triple ionization, or at very tight tolerances ($< 10^{-10}$) where depth $m = 5$ achieves a $2.9\times$ iteration reduction.

Despite the modest speedup at standard tolerances, we include Anderson acceleration in the pipeline for two reasons: (1)~it provides a convergence guarantee for the harder plasma conditions that may arise in applications beyond the parameter range tested here, and (2)~the \texttt{lax.scan} implementation adds negligible overhead when $m = 0$ (Picard fallback), so the capability is available at zero cost when not needed.

\subsection{JIT Compilation Overhead}

JAX's just-in-time compilation introduces a one-time overhead of approximately \SIrange{30}{120}{\second} for the full forward model, depending on the number of emission lines and the complexity of the computation graph.
This cost is amortized over subsequent calls: after the first \texttt{jit}-compiled execution, repeated calls incur only data transfer and kernel launch overhead.
For manifold generation workloads involving $10^4$--$10^6$ forward evaluations, the compilation overhead is negligible ($< 0.01\%$ of total runtime).
For interactive single-spectrum analysis, the compilation overhead can be mitigated by pre-compiling the model during application startup.

\subsection{GPU Memory Constraints}

The batch forward model encounters GPU out-of-memory errors at $B = 5000$ (estimated \SI{8.2}{\giga\byte})
% Source: batch_forward_results.json/results[7]/estimated_memory_gb = 8.192, gpu_oom = true
and the end-to-end pipeline at $B = 10{,}000$ (estimated \SI{91.6}{\giga\byte}).
% Source: e2e_pipeline_results.json/results[4]/gpu_error
These limits are consistent with the memory budget formula (\cref{eq:bmax}) and the \SI{32}{\giga\byte} capacity of the V100S.
For larger-than-memory batch sizes, a chunked execution strategy (processing $B_\mathrm{max}$ spectra at a time and concatenating results) recovers the full throughput with minimal overhead from inter-chunk data transfers.

\subsection{FAISS GPU Index}

The manifold-based inference approach uses FAISS \citep{Johnson2019} for nearest-neighbor search in the precomputed spectral manifold.
GPU FAISS benchmark data were not available for this study (the benchmark returned null results),
% Source: benchmark_summary.json/headline_numbers/faiss_gpu_vs_cpu_speedup_largest = null
and we do not report fabricated timing numbers.
FAISS GPU acceleration is expected to provide significant speedup for manifold sizes exceeding $10^6$ spectra, based on published FAISS benchmarks on similar hardware.
This remains an important direction for future work.

\subsection{Limitations}

Several limitations should be noted:
\begin{itemize}
  \item \textbf{Hardware specificity.} All benchmarks were performed on a single GPU architecture (Tesla V100S). Performance on other GPUs (A100, H100, consumer GPUs, Apple Silicon) may differ significantly due to differences in memory bandwidth, compute unit count, and float64 throughput. The V100S delivers float64 at half the rate of float32; on consumer GPUs where float64 throughput is $1/32$ of float32, the speedup ratios would be substantially lower.
  \item \textbf{LTE assumption.} The current implementation assumes local thermodynamic equilibrium (LTE), a single-zone plasma, and optically thin emission. Non-LTE effects, spatial inhomogeneity, and self-absorption correction are not included in the GPU-accelerated pipeline.
  \item \textbf{Single-spectrum latency.} At batch size $B = 1$, the GPU is slower than CPU ($0.21\times$ speedup for the end-to-end pipeline) due to kernel launch and data transfer overhead.
  % Source: e2e_pipeline_results.json/results[0]/speedup = 0.21
  The ``real-time'' advantage of GPU acceleration is realized only for batch sizes $B \geq 10$.
  % Source: benchmark_summary.json/headline_numbers/e2e_crossover_batch_size = 10
  \item \textbf{Float64 requirement.} The Weideman $N = 36$ coefficients span 14~orders of magnitude, and the Boltzmann plot slope determination requires high-precision arithmetic. Float32 computation would approximately double throughput but introduces $\sim 0.1$--$1\%$ errors in the Voigt profile wings and limits temperature precision to $\sim 3$~significant figures.
\end{itemize}
